%%
%% This is file `sample-sigconf.tex',
%% generated with the docstrip utility.
%%
%% The original source files were:
%%
%% samples.dtx  (with options: `all,proceedings,bibtex,sigconf')
%%
%% IMPORTANT NOTICE:
%%
%% For the copyright see the source file.
%%
%% Any modified versions of this file must be renamed
%% with new filenames distinct from sample-sigconf.tex.
%%
%% For distribution of the original source see the terms
%% for copying and modification in the file samples.dtx.
%%
%% This generated file may be distributed as long as the
%% original source files, as listed above, are part of the
%% same distribution. (The sources need not necessarily be
%% in the same archive or directory.)
%%
%%
%% Commands for TeXCount
%TC:macro \cite [option:text,text]
%TC:macro \citep [option:text,text]
%TC:macro \citet [option:text,text]
%TC:envir table 0 1
%TC:envir table* 0 1
%TC:envir tabular [ignore] word
%TC:envir displaymath 0 word
%TC:envir math 0 word
%TC:envir comment 0 0
%%
%% The first command in your LaTeX source must be the \documentclass
%% command.
%%
%% For submission and review of your manuscript please change the
%% command to \documentclass[manuscript, screen, review]{acmart}.
%%
%% When submitting camera ready or to TAPS, please change the command
%% to \documentclass[sigconf]{acmart} or whichever template is required
%% for your publication.
%%
%%
% \documentclass[sigconf]{acmart}
\documentclass[sigconf,natbib=true,anonymous=true]{acmart}

%%
%% \BibTeX command to typeset BibTeX logo in the docs
\AtBeginDocument{%
  \providecommand\BibTeX{{%
    Bib\TeX}}}

%% Rights management information. This information is sent to you
%% when you complete the rights form. These commands have SAMPLE
%% values in them; it is your responsibility as an author to replace
%% the commands and values with those provided to you when you
%% complete the rights form.
\setcopyright{acmlicensed}
\copyrightyear{2026}
\acmYear{2026}
\acmDOI{XXXXXXX.XXXXXXX}

%% These commands are for a PROCEEDINGS abstract or paper.
\acmConference[SIGIR '26]{The 49th International ACM SIGIR Conference on Research and Development in Information Retrieval}{July 20--24, 2026}{Melbourne | Naarm, Australia}
%%
%%  Uncomment \acmBooktitle if the title of the proceedings is different
%%  from ``Proceedings of ...''!
%%
%%\acmBooktitle{Woodstock '18: ACM Symposium on Neural Gaze Detection,
%%  June 03--05, 2018, Woodstock, NY}
\acmISBN{978-1-4503-XXXX-X/26/07}

%%
%% Submission ID.
%% Use this when submitting an article to a sponsored event. You'll
%% receive a unique submission ID from the organizers
%% of the event, and this ID should be used as the parameter to this command.
%%\acmSubmissionID{123-A56-BU3}

%%
%% For managing citations, it is recommended to use bibliography
%% files in BibTeX format.
%%
%% You can then either use BibTeX with the ACM-Reference-Format style,
%% or BibLaTeX with the acmnumeric or acmauthoryear sytles, that include
%% support for advanced citation of software artefact from the
%% biblatex-software package, also separately available on CTAN.
%%
%% Look at the sample-*-biblatex.tex files for templates showcasing
%% the biblatex styles.
%%
%%
%% The majority of ACM publications use numbered citations and
%% references. The command \citestyle{authoryear} switches to the
%% "author year" style.
%%
%% If you are preparing content for an event
%% sponsored by ACM SIGGRAPH, you must use the "author year" style of
%% citations and references.
%% Uncommenting
%% the next command will enable that style.
%%\citestyle{acmauthoryear}

%%
%% end of the preamble, start of the body of the document source.
\begin{document}

%%
%% The "title" command has an optional parameter,
%% allowing the author to define a "short title" to be used in page headers.
\title{Zero-Shot Multi-Agent Generation for Specialized RAG Workflows: An Empirical Evaluation}

%%
%% The "author" command and its associated commands are used to define
%% the authors and their affiliations.
%% Of note is the shared affiliation of the first two authors, and the
%% "authornote" and "authornotemark" commands
%% used to denote shared contribution to the research.
\author{First Author}
\authornote{Both authors contributed equally to this research.}
\email{author1@institution.edu}
% \orcid{1234-5678-9012}
\affiliation{%
  \institution{Institution Name}
  \city{City}
  \country{Country}
}

\author{Second Author}
\authornotemark[1]
\email{author2@institution.edu}
\affiliation{%
  \institution{Institution Name}
  \city{City}
  \country{Country}
}

%%
%% By default, the full list of authors will be used in the page
%% headers. Often, this list is too long, and will overlap
%% other information printed in the page headers. This command allows
%% the author to define a more concise list
%% of authors' names for this purpose.
\renewcommand{\shortauthors}{Author et al.}

%%
%% The abstract is a short summary of the work to be presented in the
%% article.
\begin{abstract}
Auto-generated Multi-Agent Systems (Auto-MAS) promise to automate the design of complex agentic workflows from high-level task descriptions. While effective for coding and reasoning benchmarks, their application to Retrieval-Augmented Generation (RAG) remains unexamined---a critical gap for industrial ``cold-start'' scenarios where labeled data for workflow optimization is unavailable. We hypothesize that zero-shot Auto-MAS generators systematically over-engineer retrieval workflows, introducing mechanisms such as multi-agent debates, reflexion loops, and majority voting that add cost without improving retrieval quality. To our knowledge, this is the first empirical evaluation of leading Auto-MAS frameworks on RAG benchmarks. We compare four zero-shot generators---SwarmAgentic, MetaAgent, MAS-Zero, and a structured workflow approach we propose (AutoMAW)---against manually designed baselines on HotpotQA (multi-hop retrieval) and FinanceBench (tool-augmented financial reasoning). Our results confirm the over-engineering hypothesis: the most complex generated architecture (9 agents with debate and voting) underperforms even a non-agentic baseline, while simpler declarative workflow generation achieves the highest accuracy at moderate cost. We find that structural constraints on the generation process yield more reliable retrieval pipelines than unconstrained code synthesis, and that one-time workflow generation consistently outperforms per-task generation. These findings offer practical guidance for deploying Auto-MAS in data-scarce retrieval settings.
\end{abstract}

%%
%% The code below is generated by the tool at http://dl.acm.org/ccs.cfm.
%% Please copy and paste the code instead of the example below.
%%
\begin{CCSXML}
<ccs2012>
   <concept>
       <concept_id>10002951.10003317</concept_id>
       <concept_desc>Information systems~Information retrieval</concept_desc>
       <concept_significance>500</concept_significance>
       </concept>
   <concept>
       <concept_id>10010147.10010178</concept_id>
       <concept_desc>Computing methodologies~Artificial intelligence</concept_desc>
       <concept_significance>500</concept_significance>
       </concept>
 </ccs2012>
\end{CCSXML}

\ccsdesc[500]{Information systems~Information retrieval}
\ccsdesc[500]{Computing methodologies~Artificial intelligence}

%%
%% Keywords. The author(s) should pick words that accurately describe
%% the work being presented. Separate the keywords with commas.
\keywords{automated multi-agent systems, retrieval-augmented generation, agentic workflows}

%%
%% This command processes the author and affiliation and title
%% information and builds the first part of the formatted document.
\maketitle

\section{Introduction}

Retrieval-Augmented Generation (RAG) \cite{lewis2020retrieval} is evolving from static linear pipelines\cite{hong2024metagpt, wu2024autogen, li2023camel, nguyen2025marag, maps2025} to dynamic Multi-Agent Systems (MAS) capable of planning and self-correction. However, manually engineering these workflows remains a brittle and labor-intensive bottleneck. To address this, Automated Multi-Agent Systems (Auto-MAS) have emerged, treating workflow generation as an optimization problem. Leading approaches like AFlow \citep{zhang2025aflow}, MaAS \cite{zhang2025agentic-supernet}, and GPTSwarm \cite{gptswarm24} achieve state-of-the-art results by iteratively refining architectures against labeled training sets. Similarly, MAS-GPT \cite{ye2025mas} relies on fine-tuning LLMs on extensive agent interaction corpora.

These data-dependent paradigms\cite{chen2025mmoa, asai2024selfrag} face a critical barrier in industrial applications known as the "cold-start" problem. In specialized domains such as financial forensics or proprietary tech support, labeled workflow data is unavailable, and creating it is prohibitively expensive. This necessitates zero-shot generators like SwarmAgentic \cite{zhang-etal-2025-swarmagentic}, MetaAgent\cite{qian2025metaagentselfevolvingagenttool}, MAS-Zero\cite{ke2025maszerodesigningmultiagentsystems} which synthesize agent teams from high-level descriptions without prior training. While SwarmAgentic outperforms methods like ADAS \cite{hu2024ADAS} on reasoning benchmarks, its efficacy in RAG scenarios remains unproven. Specifically, retrieval tasks introduce unique challenges such as noisy context\cite{wei2025instructrag} and strict tool orchestration that purely reasoning-based benchmarks do not capture.

This gap is critical because industrial RAG demands reliability over creativity. For instance, quantitative reasoning over financial filings (FinanceBench \cite{islam2023financebench}) requires precise coordination between retrieval and calculation tools. A hallucinated workflow that bypasses necessary tool calls is unacceptable in such high-stakes environments.

\textbf{Related Work.}
Multi-agent approaches to RAG have shown promise in controlled settings: MA-RAG~\cite{nguyen2025marag} uses plan-execute-summarize chains, MAPS~\cite{maps2025} optimizes individual RAG agents, MMOA-RAG~\cite{chen2025mmoa} applies multi-agent reinforcement learning, and MADAM-RAG~\cite{han2025madamrag} employs inter-agent debate to resolve conflicting evidence. However, all of these manually design a fixed agent topology. On the automation side, frameworks such as AFlow~\cite{zhang2025aflow}, GPTSwarm~\cite{gptswarm24}, ADAS~\cite{hu2024ADAS}, and MAS-GPT~\cite{ye2025mas} treat workflow generation as an optimization problem, achieving strong results on coding and reasoning benchmarks, but rely on iterative refinement against labeled data. Zero-shot generators like SwarmAgentic~\cite{zhang2025swarmagenticfullyautomatedagentic}, MetaAgent~\cite{qian2025metaagentselfevolvingagenttool}, and MAS-Zero~\cite{ke2025maszerodesigningmultiagentsystems} remove this data dependency, yet none have been evaluated on retrieval tasks. This paper addresses exactly this gap: we present the first empirical evaluation of zero-shot Auto-MAS generators on RAG benchmarks, testing whether architectures designed for open-ended reasoning transfer effectively to the tool-intensive, noise-sensitive demands of information retrieval.

\section{Experimental Setup}

We design our experiments to simulate a "cold-start" industrial scenario: the system receives a dataset description and a set of generic tools but has no access to training examples (Zero-Shot).

All systems share the same retrieval pipeline: BGE-M3 retrieves the top-20 results, which are then re-ranked by bge-reranker-v2-m3 to the final top-10~\cite{chen2024bge, li2023making}.

\subsection{Datasets}
We selected two benchmarks that represent contrasting dimensions of RAG complexity, ensuring a comprehensive evaluation of the generated architectures across both academic and industrial scenarios.

\textbf{HotpotQA (Multi-hop Retrieval)} \cite{yang2018hotpotqa}: An academic benchmark requiring multi-step reasoning over Wikipedia documents. We use the \textit{fullwiki} setting (open-domain retrieval) and sample a stratified subset of 500 questions from the validation set (distractor). This dataset tests the Auto-MAS's ability to autonomously discover chaining strategies (e.g., retrieving an entity, reading its content, and generating a secondary query) without explicit supervision.

\textbf{FinanceBench (Tool-Augmented Reasoning)} \cite{islam2023financebench}: An industrial benchmark consisting of 150 questions derived from real-world financial reports. Unlike HotpotQA, this task requires precise numerical extraction and calculation, stressing the system's ability to generate specialized tool-using agents (e.g., a "Calculator" or "Financial Analyst"). Crucially, it represents a data-scarce scenario where no training examples are available for workflow optimization, forcing the system to rely entirely on zero-shot architectural generation.

\subsection{Systems Under Evaluation}

We compare several approaches to multi-agent RAG, spanning non-agentic baselines, a fixed multi-agent pipeline, and four inference-time MAS generation frameworks. All systems operate in a zero-shot setting without training or validation data; for frameworks that include optimization loops in their full formulations, we evaluate only the single-pass generation variant.

\textbf{Naive RAG}. A straightforward baseline that retrieves top-k passages, concatenates them as context, and prompts the LLM to generate an answer in a single step. This non-agentic approach serves as a lower bound to assess the value of multi-agent orchestration over direct retrieval augmentation.

\textbf{Single-Agent RAG}. An iterative tool-calling agent that can perform multiple retrieval and reasoning steps before generating a final answer. The agent has access to retrieve, rerank, and calculate tools, allowing it to decompose complex questions into sub-queries, gather evidence through sequential retrieval calls, and perform numerical computations. Unlike Naive RAG's fixed pipeline, this baseline demonstrates adaptive information gathering through tool use, while still operating as a single autonomous agent without multi-agent coordination.

\textbf{MA-RAG}~\cite{nguyen2025marag} implements a plan-execute-summarize pipeline with a fixed graph topology. A planning agent decomposes the question into sub-tasks, each classified as question-answering (requiring retrieval) or aggregation (synthesizing prior results). QA steps follow a five-stage sequence---query definition, retrieval, reranking, per-document extraction, and answer generation---while aggregation steps reason over accumulated context without retrieval. A final summarization agent synthesizes all step-level outputs into the answer.

\textbf{MetaAgent}~\cite{qian2025metaagentselfevolvingagenttool} implements a learning-by-doing approach where retrieval strategies
evolve through task experience. A tool router maps natural language help requests
to retrieval tools, while self-reflection after each task refines search planning
and builds a persistent knowledge base from interaction history, enabling continuous improvement without training data. However, the system generates individual systems for each query. Hence, does not support 

\textbf{MAS-Zero}~\cite{ke2025maszerodesigningmultiagentsystems} is an inference-time framework for automatic multi-agent system design that similarly operates without requiring validation data, but focuses on designing and refining entire agent system architectures rather than retrieval strategies.  The framework uses a meta-agent to iteratively design, critique, and refine multi-agent configurations through three key steps: MAS-Init executes predefined building blocks (CoT, CoT-SC, Debate, Self-Refine), MAS-Evolve performs meta-level design by decomposing problems into sub-tasks and assigning specialized sub-MAS configurations while incorporating meta-feedback on solvability and completeness, and MAS-Verify selects the most reliable solution from all candidates.

\textbf{SwarmAgentic}~\cite{zhang2025swarmagenticfullyautomatedagentic} generates multi-agent teams from scratch given only a task description. An LLM defines a set of agent roles (with names, descriptions, and capabilities) together with a workflow specifying their collaboration protocol, then produces executable coordination code. We evaluate the single-pass generation variant, where one team configuration is produced and reused across all questions, without the iterative Particle Swarm Optimization loop described in the original work.

\textbf{AutoMAW.} In addition to the above frameworks, we propose a structured workflow generation approach (\textbf{AutoMAW}) that constructs multi-agent systems through declarative specification rather than executable code synthesis. The framework operates in two stages: first generating a pool of specialized agents with distinct roles and capabilities, then establishing coordination links between them to form a directed workflow graph. Unlike code-generating approaches, AutoMAW decouples workflow topology from runtime execution, constraining the generation space to agent roles and their coordination structure. We include this method to examine whether explicit structural constraints on the generation process offer advantages for retrieval tasks.

MetaAgent, MAS-Zero, SwarmAgentic and AutoMAW are evaluated in two configurations: \textbf{One-Time Generation} — a single MAS is generated once and reused across all benchmark tasks, and \textbf{Per-Task Generation} — a new MAS is generated independently for each task.

\subsection{Evaluation Metrics}
Given the open-ended nature of the generated answers, we use LLM-as-a-Judge~\cite{li2025genjudge} accuracy as the primary metric. GPT-4o-mini serves as the judge, comparing each system output against the gold answer and returning a binary correct/incorrect verdict. We additionally report token-level F1 on HotpotQA to enable comparison with prior work, and total inference cost (USD) per benchmark to quantify the cost--accuracy trade-off across systems.

\section{Results and Analysis}

\subsection{Main Results}

Table~\ref{tab:main_results} presents the performance of all evaluated systems across both benchmarks. All systems use GPT-4o-mini as the underlying LLM to ensure a fair comparison.

\iffalse
\begin{table}[h]
  \centering
  \caption{Accuracy and cost comparison on FinanceBench and HotpotQA. All systems use GPT-4o-mini. Best accuracy results per dataset are bolded.}
  \label{tab:main_results}
  \begin{tabular}{lccccc}
    \toprule
    & \multicolumn{2}{c}{\textbf{FinanceBench}} & \multicolumn{2}{c}{\textbf{HotpotQA}} \\
    \cmidrule(lr){2-3} \cmidrule(lr){4-5}
    \textbf{System} & Acc. (\%) & Cost (\$) & Acc. (\%) & Cost (\$) \\
    \midrule
    Naive RAG        & 43.80 & 0.20 & 56.67 & 0.11 \\
    Single-Agent     & 52.67 & 1.51 & 73.33 & 1.17 \\
    \midrule
    \multicolumn{5}{l}{\textbf{One-Time Generation}} \\
    \midrule
    SwarmAgentic     & 56.70 & ... & 71.00 & ... \\
    MetaAgent        & 52.00 & 2.41 & ... & ... \\
    MAS-Zero         & 42.90 & 0.76 & 72.60 & 0.36 \\
    AutoMAW          & ... & ... & ... & ... \\
    \midrule
    \multicolumn{5}{l}{\textbf{Per-Task Generation}} \\
    \midrule
    SwarmAgentic     & 50.00 & ... & ... & ... \\
    MetaAgent        & ... & ... & ... & ... \\
    MAS-Zero         & ... & ... & ... & ... \\
    AutoMAW          & ... & ... & ... & ... \\
    \bottomrule
  \end{tabular}
\end{table}
\fi

% \begin{table}[h]
%   \centering
%   \caption{Accuracy and context recall comparison on FinanceBench and HotpotQA. All systems use GPT-4o-mini. Best accuracy results per dataset are bolded.}
%   \label{tab:main_results}
%   \begin{tabular}{lccc}
%     \toprule
%     & \textbf{FinanceBench} & \multicolumn{2}{c}{\textbf{HotpotQA}} \\
%     \cmidrule(lr){2-2} \cmidrule(lr){3-4}
%     \textbf{System} & Acc. (\%) & Acc. (\%) & Context Recall (\%) \\
%     \midrule
%     Naive RAG        & 43.80 & 56.67 & ... \\
%     Single-Agent     & 52.67 & 73.33 & ... \\
%     MA-RAG           & ... & ... & ... \\
%     \midrule
%     \multicolumn{4}{l}{\textbf{One-Time Generation}} \\
%     \midrule
%     SwarmAgentic     & 56.70 & 71.00 & ... \\
%     MetaAgent        & 52.00 & ... & ... \\
%     MAS-Zero         & 42.90 & 72.60 & ... \\
%     AutoMAW          & ... & ... & ... \\
%     \midrule
%     \multicolumn{4}{l}{\textbf{Per-Task Generation}} \\
%     \midrule
%     SwarmAgentic     & 50.00 & ... & ... \\
%     MetaAgent        & 46.80 & ... & ... \\
%     MAS-Zero         & 32.70 & 61.40 & 70.90 \\
%     AutoMAW          & ... & ... & ... \\
%     \bottomrule
%   \end{tabular}
% \end{table}

\begin{table}[h]
  \centering
  \caption{Accuracy, F1 score, and cost comparison on FinanceBench and HotpotQA. All systems use GPT-4o-mini. Best accuracy results per dataset are bolded.}
  \label{tab:main_results}
  \begin{tabular}{lcccccc}
    \toprule
    & \multicolumn{2}{c}{\textbf{FinanceBench}} & \multicolumn{3}{c}{\textbf{HotpotQA}} \\
    \cmidrule(lr){2-3} \cmidrule(lr){4-6}
    \textbf{System} & Acc. (\%) & Cost (\$) & Acc. (\%) & F1 (\%) & Cost (\$) \\
    \midrule
    \multicolumn{6}{l}{\textit{Manually Designed Systems}} \\
    Naive RAG        & 43.80 & 0.20 & 56.67 & ... & 0.11 \\
    Single-Agent     & 52.67 & 1.51 & 73.33 & ... & 1.17 \\
    MA-RAG           & 32.00   & 0.81  & ...   & ... & ...  \\
    \midrule
    \multicolumn{6}{l}{\textit{Automated Design: One-Time Generation}} \\
    SwarmAgentic     & 56.70 & ...  & 71.00 & ... & ...  \\
    MetaAgent        & 52.00 & 2.41 & 4.30  & ... & 24.86 \\
    % MetaAgent*       & -- & -- & --  & -- & -- \\
    MAS-Zero         & 42.90 & 0.76 & 72.60 & ... & 0.36 \\
    AutoMAW          & 66.00   & 1.65  & 85.80   & 14.60 & 2.67  \\
    \midrule
    \multicolumn{6}{l}{\textit{Automated Design: Per-Task Generation}} \\
    SwarmAgentic     & 50.00 & ...  & 69.90 & ... & ...  \\
    MetaAgent        & 46.80 & 1.02 & 52.20 & ... & 1.91 \\
    MAS-Zero         & 32.70 & 0.71 & 61.40 & ... & 0.46 \\
    AutoMAW          & ...   & ...  & ...   & ... & ...  \\
    \bottomrule
  \end{tabular}

  \vspace{4pt}
\end{table}

\begin{figure*}[!h]
\centering
\begin{minipage}{0.48\textwidth}
  \centering
  \includegraphics[width=\linewidth]{pareto_hotpotqa.pdf}
  \centerline{(a) HotpotQA}
\end{minipage}\hfill
\begin{minipage}{0.48\textwidth}
  \centering
  \includegraphics[width=\linewidth]{pareto_financebench.pdf}
  \centerline{(b) FinanceBench}
\end{minipage}
\caption{Cost--accuracy Pareto fronts on both benchmarks. Systems closer to the top-left corner offer better accuracy per dollar. MetaAgent OT on HotpotQA (\$24.86, 4.3\%) is omitted for readability.}
\label{fig:pareto}
\Description{Pareto front of multi-agent system architectures on HotpotQA and FinanceBench benchmarks.}
\end{figure*}

\subsection{Discussion}

The results in Table~\ref{tab:main_results} and Figure~\ref{fig:pareto} reveal several findings that support our over-engineering hypothesis.

\textbf{Structural constraints outperform unconstrained generation.}
AutoMAW, which generates declarative workflow graphs without synthesizing executable code, achieves the highest accuracy on both benchmarks (66.0\% on FinanceBench, 85.8\% on HotpotQA), outperforming all code-generating Auto-MAS approaches. By constraining the generation space to agent roles and coordination links, it avoids the runtime failures inherent in synthesized orchestration code while retaining sufficient flexibility for multi-step retrieval.

\textbf{Complexity does not correlate with performance.}
MAS-Zero deploys the most elaborate architecture---nine agents with debate, reflexion, and majority voting---yet scores 42.9\% on FinanceBench, below even Naive RAG (43.8\%). Similarly, MA-RAG's fixed multi-agent pipeline (32.0\%) underperforms the single-step baseline despite its plan-execute-summarize design. These results indicate that adding coordination overhead to retrieval workflows without task-specific calibration degrades rather than improves accuracy.

\textbf{One-time generation is more stable than per-task.}
Across all Auto-MAS frameworks, one-time generation consistently matches or exceeds per-task generation (e.g., SwarmAgentic: 56.7\% vs.\ 50.0\% on FinanceBench; MAS-Zero: 42.9\% vs.\ 32.7\%). Regenerating a workflow for each query introduces architectural variance---different agent counts, tool bindings, and control flow---without a corresponding adaptation benefit, since no task-level feedback is available in the zero-shot setting.

\textbf{Code-generating approaches are fragile in tool-intensive RAG.}
MetaAgent achieves competitive accuracy on FinanceBench (52.0\% OT) but catastrophically fails on HotpotQA (4.3\% at \$24.86), where multi-hop queries demand precise tool orchestration across many retrieval calls. When generated code fails to correctly invoke or chain retrieval tools, the entire pipeline collapses---a failure mode absent from reasoning-only benchmarks where tool use is minimal.

\subsection{Case Study: Architecture Analysis}

To illustrate the over-engineering phenomenon, we compare the architectures generated by MAS-Zero and AutoMAW for a representative FinanceBench query: \textit{``What is the FY2018 capital expenditure amount (in USD millions) for 3M?''} Both systems produced the correct answer, allowing us to focus on architectural differences.

\textbf{MAS-Zero} (Figure~\ref{fig:casestudy}a) generates the most complex architecture among all evaluated systems: a two-phase debate-and-verify pipeline with nine agents. Retrieved documents are distributed to three parallel \textit{Debate Agents}, whose candidate answers are evaluated by a \textit{Critic Agent} performing reflexion-style assessment. If evidence is insufficient, a refined query loops back to the retriever. In the second phase, three parallel \textit{Final Agents} process re-scored evidence, and their outputs are aggregated through \textit{Majority Voting}. Each query triggers up to nine LLM calls plus potential refinement loops---substantial overhead for a straightforward numerical extraction task.

\textbf{AutoMAW} (Figure~\ref{fig:casestudy}b) produces the most minimalist design: a linear two-agent pipeline where a \textit{FinancialResearchAgent} handles retrieval and extraction, and a \textit{FinancialAnalysisAgent} performs reasoning and answer synthesis. Agent roles and coordination are defined declaratively without executable code, yielding low latency, a minimal failure surface, and an interpretable workflow.

\textbf{Comparison.} These two architectures represent opposite ends of the complexity spectrum---nine agents with debate, reflexion, and voting versus two agents in a linear chain---yet both answer correctly. The contrast directly illustrates our central finding: for tool-augmented retrieval, \textit{architectural complexity does not monotonically improve performance}. MAS-Zero's elaborate mechanisms offer self-correction potential but incur disproportionate cost (Table~\ref{tab:main_results}), while AutoMAW's constrained generation avoids the runtime failures that plague more complex synthesized workflows.

\begin{figure}[!h]
\centering
\begin{minipage}{0.48\columnwidth}
  \centering
  \includegraphics[width=\linewidth]{mas-zero_solution.pdf}
  \centerline{\small (a) MAS-Zero}
\end{minipage}\hfill
\begin{minipage}{0.48\columnwidth}
  \centering
  \includegraphics[width=\linewidth]{automaw_solution.pdf}
  \centerline{\small (b) AutoMAW}
\end{minipage}
\caption{Generated architectures for a FinanceBench query. MAS-Zero produces a 9-agent debate-and-verify pipeline; AutoMAW generates a minimal 2-agent chain. Both answer correctly.}
\label{fig:casestudy}
\Description{Comparison of MAS-Zero and AutoMAW generated architectures.}
\end{figure}


\section{Conclusion}

We presented the first empirical evaluation of zero-shot Auto-MAS generators on RAG benchmarks, comparing four frameworks across academic multi-hop retrieval and industrial financial reasoning. Our central finding is that auto-generated agents systematically over-engineer retrieval workflows: mechanisms such as multi-agent debate, reflexion loops, and majority voting---effective in reasoning-only settings---add substantial cost without improving retrieval quality, and in some cases degrade it below non-agentic baselines. We find that constraining the generation space through declarative workflow specification yields more reliable and accurate retrieval pipelines than unconstrained code synthesis, and that reusing a single generated architecture across tasks is more stable than per-task regeneration. For practitioners facing cold-start RAG deployment, our results suggest that a well-structured retrieval graph with few specialized agents is preferable to elaborate self-organizing architectures. Future work should investigate whether incorporating minimal task feedback (few-shot) can bridge the gap between zero-shot generation and fully optimized workflows.

\bibliographystyle{ACM-Reference-Format}
\bibliography{main}

\end{document}
\endinput
%%
%% End of file `sample-sigconf.tex'.

